A deterministic linear-time tape transposer, even with input-dependent
head movements, would yield binary storage-code families whose minimum
paired input--recovery distance diverges with length at each sufficiently
large fixed incidence-degree cap. Their guaranteed rate approaches one as the cap
grows. We give the quantitative dependence on the
tape count and retain the same rate-loss exponent at the machine's actual
bit positions and append times.

We then construct abstract codes satisfying a stronger rate requirement.
From a fixed-degree, odd-characteristic linear seed of rate above one half
and diverging odd girth, a tensor majority operator gives binary families
of rate at least $1-C\Delta^{-\alpha}$ at every sufficiently large degree
cap $\Delta$, for fixed $C>0$ and $0<\alpha<1/2$. The binary codes are
syndrome fibres on the original tensor coordinates. One common partition
works for every source subset, with an exponential tail for the loss in
fibre rate. Repetition and paid leftover bits extend the construction to
every sufficiently large length, preserving the degree cap and paired
separation. Thus the unrestricted slow-gap criterion for ruling out a
transposer has a false premise. Physical realization of the recoveries
remains open. We also give a three-edge obstruction for literal tensor
checks, a local-code ceiling for two-slot peeling, and information bounds
on prediction from overlapping patches of input coordinates.

All new amplification, counting and padding arguments are proved in Lean
from an explicit seed hypothesis. A written proof supplies that seed
from the earlier above-half-rate construction; its upstream arithmetic
geometry is not formalized here. The tape-transfer and prediction
theorems are also Lean-verified.
